\section{Replication Methodology}\label{replication-methodology}

This document describes the systematic methodology used for replicating
experiments from the literature.

\subsection{Definition of
``Replication''}\label{definition-of-replication}

In this repository, \textbf{replication} means:

\begin{enumerate}
\def\labelenumi{\arabic{enumi}.}
\tightlist
\item
  \textbf{Obtaining} the original source code from the paper's public
  repository (\textbf{done manually by the user} --- copied into
  \texttt{src/original/})
\item
  \textbf{Setting up} the environment (dependencies, data,
  configuration)
\item
  \textbf{Running} the original code with the original parameters
\item
  \textbf{Comparing} our output with the results reported in the paper
\item
  \textbf{Documenting} any deviations, issues, or insights
\end{enumerate}

We are \textbf{NOT} reimplementing from scratch. We are verifying that
the original code produces the claimed results.

\begin{quote}
\small \textbf{Note}: All code, dataset, and file insertion is a \textbf{manual
process} performed by the researcher. The repository scaffolding
provides the structure and conventions; the user populates each
experiment folder at their own pace.
\end{quote}

\subsection{Replication Criteria}\label{replication-criteria}

An experiment is considered \textbf{successfully replicated} when:

\begin{itemize}
\tightlist
\item
  The original code runs without errors in our environment
\item
  Key metrics (RMSE, MAE, F1, SERA, etc.) match the paper within a
  \textbf{reasonable tolerance} (±5\% for stochastic methods, exact
  match for deterministic)
\item
  Any deviations are explained (different random seeds, library
  versions, data splits)
\end{itemize}

\subsection{Standard Workflow per
Experiment}\label{standard-workflow-per-experiment}

\begin{enumerate}
\def\labelenumi{\arabic{enumi}.}
\tightlist
\item
  \textbf{Read the paper} --- understand the claims, methodology, and
  target results
\item
  \textbf{Read the ARA} (if exists) --- review compiled knowledge in the
  Obsidian Vault
\item
  \textbf{Insert the code} --- manually copy/clone the original code
  into \texttt{src/original/} (user-owned step)
\item
  \textbf{Insert datasets} --- manually place required data files
  locally (user-owned step)
\item
  \textbf{Set up environment} --- create isolated
  \texttt{venv} or \texttt{conda} env
\item
  \textbf{Run original code} --- with original parameters, document the
  process
\item
  \textbf{Record results} --- save outputs to
  \texttt{results/}, fill comparison table
\item
  \textbf{Document} --- complete README and REPLICATION\_LOG
\item
  \textbf{Update dashboard} --- run
  \texttt{scripts/update\textbackslash\{\}\_readme.py}
\end{enumerate}

\subsection{Evaluation Metrics
Reference}\label{evaluation-metrics-reference}

{\def\LTcaptype{none} % do not increment counter
\small
\begin{longtable}[]{@{}
  >{\raggedright\arraybackslash}p{(\linewidth - 4\tabcolsep) * \real{0.2759}}
  >{\raggedright\arraybackslash}p{(\linewidth - 4\tabcolsep) * \real{0.2759}}
  >{\raggedright\arraybackslash}p{(\linewidth - 4\tabcolsep) * \real{0.4483}}@{}}
\toprule\noalign{}
\begin{minipage}[b]{\linewidth}\raggedright
Metric
\end{minipage} & \begin{minipage}[b]{\linewidth}\raggedright
Domain
\end{minipage} & \begin{minipage}[b]{\linewidth}\raggedright
Description
\end{minipage} \\
\midrule\noalign{}
\endhead
\bottomrule\noalign{}
\endlastfoot
RMSE & Regression & Root Mean Squared Error \\
MAE & Regression & Mean Absolute Error \\
SERA & Imbalanced Regression & Squared Error-Relevance Area (from
IRon) \\
F1 & Classification & Harmonic mean of precision and recall \\
AUC-ROC & Classification & Area under ROC curve \\
TSS & Classification & True Skill Statistic (used in solar flare
prediction) \\
\end{longtable}
}

\subsection{Tools \& Infrastructure}\label{tools-infrastructure}

\begin{itemize}
\tightlist
\item
  \textbf{Python}: Primary language (most experiments)
\item
  \textbf{R}: Secondary language (UBA, IRon, some Torgo group papers)
\item
  \textbf{Git}: Version control for all code and results
\item
  \textbf{Obsidian Vault}: Cross-referenced knowledge base (ARAs)
\item
  \textbf{W\&B / MLflow}: Optional experiment tracking for complex runs
\end{itemize}
